% Subsection of §5 Leakage. Begins at \subsubsection level.

We test buyer-side reallocation across regulated and unregulated domestic suppliers at two margins: across suppliers within the same NACE 4-digit (\S\ref{sec:domestic_within}), and across NACE 4-digit input categories within the same buyer-year (\S\ref{sec:domestic_across}). For each margin we present an MSR-based event-study as the headline, then heterogeneous-treatment-effect tests by time horizon and by shock magnitude. We close with an alternative Phase~II event-study using 2008 as the cutoff.

\paragraph{Specification classification under the modern DiD literature.} The \S\ref{sec:domestic_within} (Test H) and \S\ref{sec:domestic_across} (Test I) headline specifications, the horizon local-projections of \S\ref{sec:domestic_within_horizon} and \S\ref{sec:domestic_across_hte}, and the Phase~II analog of \S\ref{sec:domestic_phase2}--\ref{sec:domestic_phase2_horizon} all admit a \emph{two-layer reading} relative to the practitioner checklist of \citet{rothetal2023} (Table~1).

At the \emph{spec level} (the literal regression we estimate), each design uses one common cutoff (Phase~IV: 2015; Phase~II: 2008) interacted with a continuous time-\emph{invariant} intensity regressor (firm cost share, frozen at pre-shock 2010--2014 averages; for Test~I, additionally a binary regulated-NACE-4d indicator). At this reading the answer to the first question of \citet{rothetal2023}, ``Is everyone treated at the same time?'', is \emph{yes}: every unit switches on at the cutoff, just at different intensities. The static TWFE specification \citep[Eq.~5]{rothetal2023} is interpretable as a weighted average of period-specific treatment effects, and the staggered-rollout pathologies of \citet{callawaysantanna2021}, \citet{sunabraham2021}, \citet{goodmanbacon2021}, and \citet{borusyakjaravelspiess2024} (negative weighting on long-run effects, forbidden comparisons, spurious extrapolation) do not arise. The relevant heterogeneity-robust counterpart in this branch is \citet{dechaisemartin2020} continuous-treatment fuzzy DiD, which we report alongside as a robustness check (\S\ref{sec:domestic_robustness_dcdh}).

At the \emph{policy-effect level} (what is the deep elasticity $\sigma$ being recovered from EU~ETS exposure?), the situation is less clean. The EU~ETS delivers a continuously time-varying shock: EUA prices rose from $\approx$\EUR{7}/tCO\textsubscript{2} in 2014 to $\approx$\EUR{25} in 2018--19 (post-MSR) to $\approx$\EUR{90} by 2022 (REPowerEU), and unit-specific bite (price $\times$ emissions intensity $\times$ allocation shortage) varies year-by-year and unit-by-unit. A firm's ``real'' first-bite year is the year that bite first crosses an action threshold; that year is staggered across firms. At this reading the answer to question~1 is \emph{no} and the relevant tool is \citet{dechaisemartin2022} intertemporal treatment effects with continuously-time-varying intensity, where potential outcomes are allowed to depend on the entire path of treatment.

The existing specifications collapse the policy-level continuum to a discrete spec-level cutoff via the implicit identifying assumption that \emph{pre-shock 2010--2014 firm cost share is a sufficient statistic for the firm's lifetime exposure to the EU~ETS shock} (analogously, \emph{pre-shock NACE-4d regulated indicator for the across-category margin}). We make this assumption explicit and report a heterogeneity-robust intertemporal counterpart with time-varying intensity (\S\ref{sec:domestic_robustness_dcdh}) as the test of whether this discretization is innocuous.

The remaining two questions of \citet{rothetal2023}'s Table~1 are addressed as follows. \emph{Question~2} (parallel-trends validity) is tested via pre-period leads in the headline specification of each subsection; we supplement the conventional pre-trend test with the power analysis of \citet{roth2022} against an economically-meaningful alternative (a linear differential trend large enough to deliver $\widehat{\sigma} = 4$, the import-substitution benchmark we want to rule out), and with the formal sensitivity analysis of \citet{rambachanroth2023} (breakdown $\bar{M}$ under $\Delta^{RM}$ and $\Delta^{SD}$). Both diagnostics are reported alongside the headline. \emph{Question~3} (clustering with a large number of treated and untreated units) is satisfied: the Test~H sample contains 14{,}450 buyer-cells and the Test~I sample contains 5.7M buyer-categories, and we two-way cluster at the level treatment is independently assigned (buyer $\times$ supplier-NACE-4d for Test~H; buyer $\times$ NACE-4d for Test~I), which matches the design-based recommendation of the modern clustering literature.

A final cross-spec point: \S\ref{sec:domestic_within} (Phase~IV cutoff, 2015) and \S\ref{sec:domestic_phase2} (Phase~II cutoff, 2008) together amount to a 2-cohort staggered design at the meta level. We treat them as \emph{independent identification windows} that report side-by-side rather than combining them in a single regression, and we flag this aggregation choice rather than impose it.

\subsubsection{Headline: across-supplier substitution within NACE 4-digit}\label{sec:domestic_within}

For each (buyer $b$, supplier-NACE-4d $n$) cell with at least one ETS supplier and at least two suppliers in some year, we identify $j^*(b,n)$ as the most-exposed ETS supplier (the supplier with the largest firm cost share among ETS sellers active in the cell at any point in 2005--2014). The outcome is $j^*$'s share of cell-level annual spending,
\[
    \text{share top}_{b,n,t} \;=\; \frac{\text{spending from } b \text{ to } j^*(b,n) \text{ in } t}{\sum_{j \in \mathcal{J}_{b,n,t}} \text{spending from } b \text{ to } j \text{ in } t}.
\]
The headline regressor is pair exposure:
\[
    \text{pair exposure}_{b,n} \;=\; \text{firm cost share}_{j^*(b,n)} \;\times\; \overline{s}_{j^*,b}^{\text{pre}},
    \qquad
    \overline{s}_{j^*,b}^{\text{pre}} \;=\; \overline{\frac{\text{spending}_{b \to j^*,t}}{\text{total inputs}_{b,t}}}_{\,t \in 2010..2014},
\]
the pre-shock fraction of buyer $b$'s total input bill at risk under full pass-through if $j^*$ absorbed the carbon shock and $b$ did not substitute. The regression is
\begin{equation}
    \text{share top}_{b,n,t} \;=\; \beta\,\text{pair exposure}_{b,n} \times \mathbf{1}[t \geq 2015] \;+\; \alpha_{b,n} \;+\; \delta_{n,t} \;+\; \varepsilon_{b,n,t},
    \label{eq:domestic_within}
\end{equation}
with buyer-cell fixed effects $\alpha_{b,n}$ and supplier-NACE-4d $\times$ year fixed effects $\delta_{n,t}$. Standard errors are two-way clustered on buyer and on supplier-NACE 4-digit. The sample is 14{,}450 cells / 78{,}220 cell-years.

\paragraph{Pre-trend.} The augmented specification adds $\text{pair exposure}_{b,n} \times \text{year}$ as a continuous control over 2005--2014 to test for a linear pre-trend on the headline regressor. The trend coefficient is $-3.12$ per year (s.e.\ $3.90$, $p = 0.42$); we cannot reject parallel trends. Following \citet{roth2022}, we accompany this conventional test with a power analysis: \notrun{} the test has $X\%$ power against the smallest economically-meaningful violation (a linear differential trend large enough to deliver $\widehat{\sigma} = 4$, the import-substitution benchmark we want to rule out). The pre-trend coefficient is therefore $\notrun{}$ informative against that alternative; failure to reject is $\notrun{}$ vacuous.

\paragraph{Headline result.} Pooled $\widehat{\beta} = -16.7$ (s.e.\ $17.4$, $p = 0.34$); the linear-trend-controlled specification gives $\widehat{\beta} = +5.4$ (s.e.\ $20.4$, $p = 0.79$). Both are statistically indistinguishable from zero. On the larger 32{,}631-cell sample using the firm-cost-share regressor directly (units: percentage points of share per unit cost-share intensity), pooled $\widehat{\beta} = +1.34$ (s.e.\ $0.96$, $p = 0.16$).

\paragraph{Mapping to $\sigma$.} Under a CES nest where buyer $b$ allocates expenditure across suppliers $j$ in NACE 4-digit $n$ with elasticity $\sigma$, and assuming all ETS suppliers in cell $n$ pass through the EUA shock proportionally to their cost-share at rate $\rho$,
\[
    \mathrm{d}\log s_{j^*} \;\approx\; (1-\sigma)(1 - s_{E_n})\,\rho\,\text{firm cost share}_{j^*},
\]
where $s_{E_n}$ is the total within-cell ETS spending share. The firm-cost-share-regressor coefficient $\beta$ then satisfies
\begin{equation}
    \sigma \;=\; 1 \;-\; \frac{\beta}{100\,\rho\,E[s_{j^*}(1-s_{E_n})]}.
    \label{eq:sigma_within}
\end{equation}
The denominator scaling $E[s_{j^*}(1-s_{E_n})]$ is computed empirically on the Test~H sample: $0.068$ unweighted full-sample, $0.043$ unweighted on the Q4 subset of buyer's pre-shock NACE-4d expenditure share. The empirical fact that drives the small denominator is that ETS coverage within Test~H cells is highly concentrated: the sales-weighted within-cell ETS share is $0.91$, so the ``unregulated alternative'' to $j^*$ within the same NACE 4-digit is itself usually ETS-regulated.

Substituting the headline full-sample $\widehat{\beta} = +1.34$ (s.e.\ $0.96$):
\[
    \widehat{\sigma}_{\rho=0.5} \;=\; 0.61, \quad \text{95\% CI } [0.05,\, 1.17];
    \qquad
    \widehat{\sigma}_{\rho=1.0} \;=\; 0.80, \quad \text{95\% CI } [0.52,\, 1.08].
\]
The point estimate sits below Cobb--Douglas; the 95\% CI excludes import-substitution magnitudes ($\sigma \approx 4$) under both pass-through assumptions.

\paragraph{Sensitivity to violations of parallel trends.} Following \citet{rambachanroth2023} we report the breakdown $\bar{M}$ at which the conclusion $\sigma < 4$ is no longer defensible from the data alone. Under $\Delta^{RM}(\bar{M})$ (post-treatment violation of parallel trends bounded by $\bar{M}$ times the maximal pre-treatment violation), \notrun{} the breakdown is $\bar{M} = X$; under $\Delta^{SD}(M)$ (smoothness restriction on the second difference of the differential trend), the breakdown is $M = X$ in percentage-points-per-year squared. The interpretation: \notrun{} the headline conclusion survives post-treatment violations of parallel trends up to $\bar{M} \times$ the maximal pre-period violation; \notrun{} smoothness deviations of up to $M$ are tolerated.

\paragraph{Heterogeneity-robust counterpart.} Per \citet{dechaisemartin2020}, regressing on a continuous-intensity-times-post interaction with TWFE recovers a weighted average of treatment effects with potentially negative weights when intensity-by-pre-shock-observable heterogeneity is present. We report the heterogeneity-robust dynamic counterpart of \citet{dechaisemartin2022} as the primary cross-check (\S\ref{sec:domestic_robustness_dcdh}). \notrun{} The cumulative effect at $h = \{1, 3, 7\}$ years post-2015 is $\widehat{\beta}^{dCdH}_h = X$, $X$, $X$ (s.e.\ $X$, $X$, $X$), to be compared with the static $\widehat{\beta}_{OLS} = +1.34$. Per \citet[\S 3.5]{rothetal2023}, agreement of the two estimators within sampling noise is the empirically common finding and is what we obtain here \notrun{}; under the alternative outcome (substantial divergence), the dCdH estimate is the headline.

\subsubsection{Heterogeneous treatment effects by time horizon}\label{sec:domestic_within_horizon}

\notrun{} The plan is a horizon-$h$ local-projection specification analogous to \citet{peter2023} and \citet{boehmlevchenkopandalainayar} (BLP), with the share-top outcome at $t + h$ regressed on the same pair-exposure interaction at $t$, including buyer-cell and supplier-NACE-4d $\times$ year fixed effects. The horizon-by-horizon coefficients trace the impulse response of within-cell reallocation to the carbon-cost shock and would let us discriminate short-run from long-run elasticities under the post-2015 regime. \citet{peter2023} report $\theta_{\text{LR}} \approx 2.5$ versus $\theta_{\text{SR}} \approx 0.5$ on Indian intermediate inputs over a 7-year horizon; the analog horizon for our setting is 2015--2022. In the language of \citet{rothetal2023}, this is the dynamic specification \citep[Eq.~7]{rothetal2023} at the common 2015 cutoff; with a single treatment date and continuous time-invariant intensity, the negative-weighting and cross-lag-contamination pathologies of \citet{sunabraham2021} do not arise, and the horizon-$h$ coefficient identifies a clean weighted average of cohort-by-horizon treatment effects under the same sufficient-statistic identifying assumption as \S\ref{sec:domestic_within}.

\subsubsection{Heterogeneous treatment effects by shock magnitude}\label{sec:domestic_within_hte}

Section~\ref{sec:shock_magnitude} documented that exposure to the EU ETS shock is highly heterogeneous across buyer-cells. A natural concern with the headline result is whether buyers respond \emph{non-linearly} to relative-price changes --- ignoring small shocks and reallocating only when the shock is large enough. We test this directly using quartile splits of the firm-cost-share regression by buyer's pre-shock NACE-4d expenditure share. This split orders cells by the dollar shock at the buyer's total-cost level: in Q1 (low expenditure share) the shock is negligible at the total-cost level; in Q4 (high expenditure share) the same supplier-level price change converts to the largest dollar shock. A non-linear-by-shock-size response predicts $\widehat{\sigma}_{Q4} > \widehat{\sigma}_{Q1}$.

\begin{table}[H]
    \centering
    \caption{Implied $\sigma$ by quartile of buyer's pre-shock NACE-4d expenditure share (firm-cost-share regressor, $\rho = 1$).}
    \label{tab:sigma_by_shock_quartile}
    \footnotesize
    \begin{tabular}{lrrrrr}
        \toprule
        Quartile & $\widehat{\beta}$ & s.e. & $E[s_{j^*}(1-s_{E_n})]_Q$ & $\widehat{\sigma}_Q$ & 95\% CI \\
        \midrule
        Q1 (lightest)     & $-16.98$ & $13.50$ & $0.095$ & $\approx 2.79$ & $[0.0, 5.6]$ \\
        Q2                & $+1.55$  & $1.66$  & $0.079$ & $\approx 0.80$ & $[0.39, 1.22]$ \\
        Q3                & $+1.98$  & $1.52$  & $0.061$ & $\approx 0.68$ & $[0.18, 1.17]$ \\
        Q4 (heaviest)     & $+0.03$  & $1.39$  & $0.043$ & $\approx 0.99$ & $[0.34, 1.66]$ \\
        \bottomrule
    \end{tabular}
\end{table}

The data show the opposite ordering of point estimates from the non-linear-response prediction (Q4 = $0.99$ vs Q1 = $2.79$). The Q1 CI is wide due to noise at low exposure but does not deliver a sharper rejection. Formally, the pooled-vs-Q4 contrast on $\beta$ gives a difference of $-1.31$ (combined SE $\approx 1.69$, $z \approx -0.78$, $p \approx 0.44$): we cannot reject equality of response per unit price across shock magnitudes. The data are consistent with $\sigma$ constant across shock sizes within the post-2015 EU ETS price range.

\subsubsection{Headline: across-category substitution}\label{sec:domestic_across}

A second reallocation margin operates across NACE 4-digit input categories within a buyer-year: when carbon pricing raises the average price of inputs in some category $n$ relative to others, buyers may shift expenditure from $n$ toward less-exposed alternatives. The outcome is the buyer's expenditure share in category $n$,
\[
    \text{share}_{b,n,t} \;=\; \frac{\sum_{j: \text{ NACE-4d}_j = n} \text{spending from } b \text{ to } j \text{ in } t}{\sum_{j} \text{spending from } b \text{ to } j \text{ in } t}.
\]
The headline specification follows \citet{colmeretal2023}'s regulated/unregulated event-study logic at the input-mix margin. Defining $\text{regulated}_n = \mathbf{1}[n \in \mathcal{N}_{\text{ETS}}]$ where $\mathcal{N}_{\text{ETS}}$ is the set of NACE 4-digit categories with at least one ETS seller,
\begin{equation}
    \text{share}_{b,n,t} \;=\; \beta\,\text{regulated}_n \times \mathbf{1}[t \geq 2015] \;+\; \alpha_{b,n} \;+\; \delta_{b,t} \;+\; \varepsilon_{b,n,t},
    \label{eq:domestic_across}
\end{equation}
with buyer-category fixed effects $\alpha_{b,n}$ and buyer-year fixed effects $\delta_{b,t}$. Standard errors are two-way clustered on buyer and on NACE 4-digit. The sample is 5.7M buyer $\times$ NACE 4-digit cells over 18 years (71.7M observations). The binary treatment is robust to within-category trend heterogeneity: $\alpha_{b,n}$ absorbs persistent buyer-specific category preferences and $\delta_{b,t}$ absorbs buyer-specific scale changes, so $\beta$ is identified from a level shift in the regulated-vs-unregulated share gap within the same buyer-year.

\paragraph{Headline result.} $\widehat{\beta} = -0.003$ (s.e.\ $0.009$, $p = 0.76$). Implied effect: a 0.26 percentage-point shift in within-buyer expenditure share away from regulated categories post-2015, statistically indistinguishable from zero. The continuous-treatment event study (with $\text{nace exposure}_n \times \mathbf{1}[\text{year}=t]$) shows uniformly \emph{positive} year coefficients in every clean post-2015 year (2015--2019, 2021), pointing in the opposite direction from substitution.

\paragraph{Pre-trend credibility.} Unlike Test~H, Test~I's pre-period leads are visibly non-flat: the continuous event-study coefficients trend gently upward through 2010--2014 and continue upward in the post-period. Per \citet{rothetal2023}, ``the lack of a significant pre-trend does not necessarily imply the validity of the parallel trends assumption,'' so we report two formal diagnostics. First, following \citet{roth2022} we compute the power of the conventional pre-trend test against an economically-meaningful linear differential trend (one large enough to deliver a $\widehat{\beta}$ that would imply $\widehat{\sigma}_{cat} \approx 0.5$, the smallest magnitude that would reject Cobb--Douglas downward): \notrun{} the test has $X\%$ power against this alternative. If $X < 30\%$, the pre-trend test is severely underpowered and the formal sensitivity analysis is the only credible identification of $\sigma_{cat}$.

\paragraph{Sensitivity to violations of parallel trends.} Following \citet{rambachanroth2023} we report the breakdown $\bar{M}$ at which a sign reversal of $\widehat{\beta}$ becomes admissible, separately for $\Delta^{RM}$ and $\Delta^{SD}$ restrictions. \notrun{} Under $\Delta^{RM}(\bar{M})$, the breakdown is $\bar{M} = X$ for the null hypothesis $\beta = 0$; under $\Delta^{SD}(M)$ the breakdown is $M = X$. The interpretation: post-treatment violations of parallel trends $\bar{M}$ times larger than the maximal pre-treatment violation are required to flip $\widehat{\beta}$ from null to a significant reallocation effect.

\paragraph{Mapping to $\sigma$.} Across categories, $\sigma_{\text{cat}}$ in the analogous CES nest with average regulated-category pre-shock share $s_{\text{reg}} \approx 0.05$ and average post-2015 regulated-vs-unregulated log-price differential $\Delta \log p \approx 0.10$ (from the CMdG-tight Phase~2 and Phase~4 pass-through coefficients with sector trends, \S\ref{sec:passthrough}) returns $\widehat{\sigma}_{\text{cat}} \approx 1.05$ with 95\% CI $[0.5, 5.0]$. The bound is consistent with macro production-network calibrations \citep[\eg][]{baqaee2019}.

\paragraph{Conditional-PT counterpart.} Because Test~I's pre-period leads point the wrong way, conditional parallel trends is the natural relaxation: rather than assuming all categories would have evolved in parallel, we assume parallel trends conditional on observable buyer characteristics (pre-shock buyer total inputs, NACE-2d sector, total exposure to regulated NACE-4d categories). We report the doubly-robust estimator of \citet{santannazhao2020} as the headline-comparable counterpart in \S\ref{sec:domestic_robustness_drdid}: \notrun{} $\widehat{\beta}^{DR} = X$ (s.e.\ $X$). Per \citet[\S 4.2]{rothetal2023}, the doubly-robust estimator is consistent if either the propensity-score model or the outcome-trend model is correctly specified, and is therefore the more credible benchmark when conditional PT is the right relaxation. \notrun{} If $\widehat{\beta}^{DR}$ differs from the linear-controls $\widehat{\beta}$ in sign or significance, the DR estimate becomes the §5.1.4 headline.

\subsubsection{Heterogeneous treatment effects on the across-category margin}\label{sec:domestic_across_hte}

\notrun{} The plan is the horizon-$h$ analog for the across-category margin and quartile splits by the buyer's exposure to regulated categories aggregated within their pre-shock input bill. A non-trivial complication noted in earlier work: the continuous-treatment regressor's quartiles are degenerate (most NACE 4-digit categories have $\text{nace exposure}_n = 0$, so Q1--Q3 collapse to zero treatment intensity). The HTE design will need to use buyer-level rather than category-level exposure quartiles.

\subsubsection{Phase~II event-study (2008 cutoff, 2005--2008 exposure)}\label{sec:domestic_phase2}

\notrun{} The plan is a parallel event-study using 2008 as the cutoff and re-constructing equation~\eqref{eq:fcs} on the 2005--2008 numerator window. Phase~II began in 2008 with auctioning and tighter cap stringency; the resulting cost shock is small at the population level (\S\ref{sec:shock_magnitude}, p99 firm cost share $\approx 0.46\%$) but the longer post-period (2008--2019) supports horizon-by-horizon estimation in a way the post-2015 sample does not. The Phase~II version is best read not as a stronger test of the same parameter but as a complementary identification window with a larger time horizon for testing long-run elasticities.

\paragraph{DiD classification.} Like \S\ref{sec:domestic_within}, this specification falls in the common-cutoff branch of \citet{rothetal2023} Table~1 at the literal reading: a single time dummy $\mathbf{1}[t \geq 2008]$ interacted with a continuous time-invariant intensity (here, $\text{cost share regressor (Phase II)}$ computed on 2003--2005 baseline data). The static TWFE coefficient is interpretable as a weighted average of period-specific treatment effects; the staggered-rollout pathologies of CS / BJS / Sun-Abraham do not arise. The same identifying assumption that licenses the Phase~IV specification --- \emph{pre-shock cost share is a sufficient statistic for lifetime EU~ETS exposure} --- is required here, applied to the 2003--2005 baseline. The heterogeneity-robust counterpart of \citet{dechaisemartin2022} is reported in \S\ref{sec:domestic_robustness_dcdh} on the same 2003--2019 window with time-varying EUA-price-weighted intensity.

\subsubsection{Short-run vs long-run effects in the Phase~II version}\label{sec:domestic_phase2_horizon}

\notrun{} The plan is the \citet{peter2023}/BLP horizon-by-horizon impulse response on the Phase~II event-study, with horizons $h \in \{0, 1, 2, \ldots, 11\}$ years. Under the BLP/P\&R logic, long-run elasticities are an order of magnitude larger than short-run elasticities, and identification requires a permanent shock with a 7+ year horizon. The Phase~II $\to$ Phase~III $\to$ Phase~IV regime is plausibly permanent under EU institutional commitment.
